%------------------------------------------------------------------------------%
\begin{question}
  Determine a posição relativa dos planos
  \[
    \pi_{1}\colon 2x - y + 3z - 4 = 0
  \]
  \[
    \pi_{2}\colon 4x - 2y + 6z + 5 = 0
  \]
\end{question}

%------------------------------------------------------------------------------%
\begin{resolution}{Solução}
  Os vetores normais são
  \[
    \vec{n}_{1} = \vetortri{2}{-1}{3}
    \qquad\qquad
    \vec{n}_{2} = \vetortri{4}{-2}{6}
  \]

  \pause
  Como
  \[
    \vec{n}_{2} = 2\vec{n}_{1}
  \]
  \pause
  os planos são paralelos ou coincidentes
\end{resolution}

%------------------------------------------------------------------------------%
\begin{resolution}{Solução}
  Multiplicarmos a primeira equação por 2,
  \pause
  obtemos
  \[
    4x-2y+6z-8=0
  \]
  \pause
  A segunda equação é
  \[
    4x-2y+6z+5=0
  \]
  \pause
  As constantes são diferentes

  \pause
  Portanto, as equações não representam o mesmo plano

  \pause
  Os planos são \emph{paralelos distintos}
\end{resolution}

%------------------------------------------------------------------------------%
\endinput
